\section{A Small Hub at Mid-Century}

The Depression-era parish is visible first through fragments. The 1929 financial-statement title suggests that money and institutional status were already being formally assessed as the national economy contracted. The document itself was not examined, so it cannot support a claim about debt or solvency. Civic population stood at 1,424 in 1930, slightly below the 1920 count. Yet Saint Anthony remained a base for a territory much larger than the city.

Annual directories make that continuity legible through personnel snapshots. Finbarr O'Callaghan appears in 1935 with Covelo attached. Celestine Quinlan is listed in the 1940s, together with the wider reservation-station field. Donatus Ahern appears in 1945. Isidore Kennedy's Capuchin authority record places him as Willits pastor in the mid-1950s; Alban Cullen appears in 1960. These names establish presence in particular years, not a complete uninterrupted sequence. Directory publication lag, temporary administration, and intra-order transfers make a pastor table look more exact than the underlying evidence permits.

The repeated ``(1923)'' after the parish name is more stable. By the 1930s and through the mid-century editions checked here, the national directory treated 1923 as Saint Anthony's conventional establishment date. That consistency is why the modern return to 1903 deserves notice rather than silent acceptance. Institutional memory had long attached its identity to the resident Capuchin turn.

The town grew after the Second World War. Census counts rose from 1,625 in 1940 to 2,091 in 1950 and 3,410 in 1960. Timber, transport, public services, and regional commerce formed the civic setting, but no annual Catholic census or sacramental series was located. Population growth cannot be converted into parish growth. It increased the number of possible constituencies while automobiles could also carry Catholics farther and change the practical meaning of parish territory.

The building record for these decades is especially thin. No controlled public evidence located here dates a major Willits church replacement, expansion, school, hall, or postwar remodeling. The archival photograph and Quinn attribution secure the early church story only in outline. Current appearance cannot fill the mid-century gap. County assessor records, building permits, fire-insurance maps, parish ledgers, newspapers, and Capuchin correspondence remain essential before a physical chronology can be claimed.

What the available record does establish is a durable scale of operation. Saint Anthony was large enough to remain an independently listed parish and small enough to be repeatedly staffed within the Capuchin county network. It functioned as a hub: local Sunday church, clerical residence, and departure point for remote stations. Its history at mid-century lies less in a single celebrated project than in the repeated work of keeping that system operating.

\statusframe{Directory snapshots from 1935 through 1960; federal census counts; Capuchin personnel descriptions.}{Saint Anthony remained a Capuchin-served parish conventionally dated to 1923, maintained the Covelo and at times wider station field, and served a growing small city.}{The snapshots are not a complete clergy roster, membership series, sacramental history, financial audit, or building chronology.}
